\documentclass[12pt]{article}

% ============================================
% PACKAGES
% ============================================
\usepackage{xurl}
\usepackage{url}

% Page layout
\usepackage[margin=1in]{geometry}

% Font and encoding
\usepackage[T1]{fontenc}
\usepackage[utf8]{inputenc}
\usepackage{newtxtext, newtxmath}
\usepackage{indentfirst}
\usepackage{tcolorbox}

% Math
\usepackage{amsmath}
\usepackage{amssymb}

% Graphics and figures
\usepackage{graphicx}
\usepackage{float}

% Caption formatting (APA 7: bold label, line break, italic text, left-aligned)
% Must load before subcaption
\usepackage[labelfont=bf,
            labelsep=newline,
            textfont=it,
            justification=raggedright,
            singlelinecheck=false]{caption}
\usepackage{subcaption}

% Tables
\usepackage{booktabs}
\usepackage{tabularx}
\usepackage{multirow}
\usepackage{array}

% Citations and bibliography
\usepackage{natbib}
\bibliographystyle{apalike} 

% Hyperlinks (load last among most packages)
\usepackage{hyperref}
\hypersetup{
    colorlinks=true,
    linkcolor=black,
    citecolor=black,
    urlcolor=blue,
    breaklinks=true
}

% Line spacing
\usepackage{setspace}
\onehalfspacing

% For appendix
\usepackage{appendix}

% For landscape pages (if needed for wide tables)
\usepackage{pdflscape}

% For editorial comments and TODOs
\usepackage{xcolor}
\newcommand{\new}[1]{\textcolor{blue}{#1}}  % For marking new content added in revision

% ============================================
% CUSTOM COMMANDS
% ============================================

% For easier table notes
\newcommand{\tablenote}[1]{\par\smallskip\footnotesize #1}

% ============================================
% TITLE PAGE SETUP
% ============================================

\title{\Large\textbf{The Impact of Lottery Incentives:\\ 
Increasing Response Rates Can Decrease Survey Accuracy}}

\author{
    Jerold Pearson\thanks{Email: jpearson@alumni.stanford.edu} 
    \and 
    Roger E. Levine\thanks{x}
    \and
    Jon A. Krosnick\thanks{Stanford University. Email: krosnick@stanford.edu}
    \and
    Resty Fufunan\thanks{Stanford Institute for Research in the Social Sciences. Email: resty@stanford.edu}
    \and
    Josearmando Torres\thanks{Stanford Institute for Research in the Social Sciences. Email: jtorres0@stanford.edu}
}

\date{\today}

% ============================================
% DOCUMENT BEGINS
% ============================================

\begin{document}

% Title
\maketitle

% ============================================
% ABSTRACT
% ============================================

\begin{abstract}
\noindent
When offering incentives with a fixed budget, survey researchers can conduct lotteries with many winners of smaller prizes or with fewer winners of larger prizes. The expected value of the reward for a respondent is the same, but one approach may be more psychologically compelling than the other. In an online survey of a listed sample, the present study randomly assigned each potential respondent to be offered five chances to win \$100, one chance to win \$500, or no incentive. Lottery incentives produced a marginal increase in overall response rates (32.5\% vs.\ 30.6\%, p = .055), and the two prize structures performed identically. Incentivized respondents completed the survey more thoroughly and their responses were statistically indistinguishable from those of non-incentivized respondents. More importantly, incentives disproportionately increased participation among subgroups already over-represented in the no-incentive condition---particularly women and donors---widening rather than narrowing existing participation gaps. Researchers should assess baseline heterogeneity in response propensity before deploying lottery incentives, as such strategies may increase overall response while simultaneously amplifying nonresponse bias.

\end{abstract}

\thispagestyle{empty}  % No page number on title page
\newpage
\setcounter{page}{1}

% ============================================
% INTRODUCTION
% ============================================

\section{Introduction}

Survey researchers have long been attentive to response rates (Steeh, 1981), treating them as an important indicator of survey quality (Backstrom \& Hursh, 1963; Babbie, 1990; Aday, 1996; Rea \& Parker, 1997). But the more fundamental concern is representativeness. Certain groups tend to respond at lower rates: men refuse more often than women, and lower-income households are less likely to participate (Groves, 2006; Bee \& Rothbaum, 2024). Other attributes such as socioeconomic status and racial minority status show variable effects depending on survey mode, topic, and sponsor (Groves, 2006). When these participation gaps correlate with the variables a survey aims to measure, nonresponse bias compromises accuracy regardless of the overall response rate.

This problem has become more pressing as response rates have declined across all modes: mail (CASRO, 1982), telephone (Curtin, Presser, \& Singer, 2005), and online (Sheehan, 2001). Researchers have experimented with survey length, advance notifications, personalization, follow-up protocols, and refusal conversion strategies (Brick \& Williams, 2013). Among these approaches, incentives remain one of the most studied and debated tools for increasing participation (Singer \& Ye, 2013).

Two theoretical frameworks help explain why incentives might motivate survey participation. The first is Cialdini's (1984) principle of reciprocity: when someone gives us something, we feel psychologically obligated to give something back. This obligation can be triggered by gifts, favors, or perceived concessions, and it operates even when the initial offering was unsolicited. In the survey context, an incentive represents a gift from the sponsor, creating a sense of debt that respondents may settle through participation.

The second framework is leverage-saliency theory (Groves, Singer, \& Corning, 2000). This theory proposes that different attributes of a survey request carry different "leverage" for different people. Whether an attribute influences behavior depends on whether it becomes salient during the request. Incentives, on this account, work by increasing the salience of potential rewards for individuals who assign positive leverage to monetary gain. This explains the inconsistent findings in the incentives literature: effectiveness depends on both individual predispositions and what the survey request emphasizes.

Both frameworks suggest that incentives should increase participation among people who would otherwise lack sufficient motivation to respond. If true, incentives could improve sample representativeness by drawing in underrepresented groups. But when incentives take the form of lotteries, additional psychological factors come into play.

Lottery incentives are appealing to researchers because they offer the possibility of meaningful prizes at low per-respondent cost. But lotteries introduce psychological complexity. When a researcher has a fixed budget, they can structure a lottery with many winners of smaller prizes or fewer winners of larger prizes. The expected value can be held constant, but the two framings may differ in their psychological appeal.

Kahneman and Tversky's (1979) prospect theory offers insight into how people evaluate such gambles. Two findings are relevant. First, people overweight small probabilities, which helps explain the general appeal of lotteries: a tiny chance of a gain receives more psychological weight than its objective probability warrants. Second, the value function for gains is concave, meaning that each additional dollar matters less than the one before it. A \$500 prize is not experienced as five times more valuable than a \$100 prize.

Applied to survey lotteries, these principles suggest that five chances to win \$100 may be more motivating than one chance to win \$500. Five chances feels like better odds, and the concave value function compresses the perceived difference between the two prize amounts. The alternative prediction is that prize magnitude dominates: \$500 is a more attention-grabbing number, and respondents may focus on what they could win rather than their chances of winning. A third possibility is that respondents approximate expected value despite uncertainty, in which case the two framings should perform equally. This calculation is complicated by the fact that respondents do not know how many others will enter the lottery.

Prior studies of lottery framing in surveys have produced mixed results. Surveys conducted within institutional populations present a particular puzzle for incentive research. People with strong ties to the sponsoring institution are likely already motivated to participate. Donors, association members, and those with deeper institutional involvement may respond at high rates without any incentive. For these individuals, incentives may add little.

By contrast, those with weaker institutional ties may lack intrinsic motivation to respond. According to leverage-saliency theory, incentives should matter most for this group, because the monetary offer provides leverage that the sponsor relationship does not. If this is correct, incentives would improve sample representativeness by increasing participation among the weakly affiliated.

But there is an alternative possibility. Cialdini's reciprocity principle suggests that incentives might actually amplify participation among the already-engaged. Donors, for instance, have given to the institution. When the institution offers something in return, donors may feel the reciprocity norm particularly strongly. If so, incentives could increase response most among those who already respond at high rates, worsening rather than improving representativeness.

This study examines lottery incentives in an online survey of Stanford University alumni. We randomly assigned potential respondents to one of three conditions: no incentive, five chances to win \$100, or one chance to win \$500. This design allows us to test three questions: whether lottery incentives increase response rates relative to no incentive, whether the framing of the lottery affects response, and whether incentive effects vary across subgroups defined by institutional affiliation.

We test the following hypotheses:
\begin{itemize}
    \item Hypothesis 1: Lottery incentives will increase overall response rates compared to no incentive.
    \item Hypothesis 2: The framing with more chances to win (5 × \$100) will produce higher response rates than the framing with a larger single prize (1 × \$500).
    \item Hypothesis 3: Incentive effects will be larger among weakly affiliated subgroups (non-donors, non-members, graduate-only alumni), thereby improving sample representativeness.
\end{itemize}

These hypotheses follow logically from theory and prior research. Whether they hold in practice is an empirical question with implications for how survey researchers should deploy limited incentive budgets.

\section{Literature Review}

The existing literature on lottery incentives presents a puzzle. Lotteries offer a cost-effective alternative to guaranteed payments, yet experimental studies reveal substantial heterogeneity in their effectiveness. Some find large positive effects; others find nothing at all. This review examines the sources of this heterogeneity, with particular attention to the understudied question of whether lottery incentives affect not just response rates but also the composition of the responding sample.

\subsection{Population Characteristics and Lottery Effectiveness}

The most consistent pattern in the literature is that lottery-effectiveness varies by population. Student and younger populations show reliable positive effects. Laguilles et al. (2011) and Larose et al. (2014) found significant increases among college students, as did Sanchez-Fernandez et al. (2010) and Porter and Whitcomb (2003). Martinson et al. (2000) similarly found that younger respondents were more responsive to lottery incentives.

General population surveys also tend to show positive effects. Olsen et al. (2012), Johansson and Berglund (1997), Pedersen and Nielsen (2016), Woodward (1985), and Marrett et al. (1992) all documented significant increases in response rates. Early studies established the basic finding: McCool (1991) found a 12 percentage point increase with a \$50 savings bond lottery, Blythe (1986) observed an 11 percentage point gain among social workers, and Balakrishnan et al. (1992) reported a 6.8 point increase in a large consumer survey. International replications suggest these effects generalize beyond U.S. samples (Angur et al. 1995; Johansson and Berglund 1997; Conn et al. 2024).

The pattern reverses for high-income, highly educated, and professional populations. Null findings are common in surveys of physicians and healthcare professionals (Moses et al. 2004; Cottrell et al. 2015; Grava-Gubins and Scott 2008; Halpern et al. 2011; Ulrich et al. 2005; Dykema et al. 2011; Gallagher et al. 2001; McLaren and Shelley 2000; So et al. 2018). Surveys of older adults show similar null results (Finsen et al. 2006; Gjøstein et al. 2016; Roberts et al. 2000), as do surveys of professionals and executives (Lorenzi et al. 1988; Coryn et al. 2020; Golden et al. 1981). Warriner et al. (1996), Harris et al. (2008), Lipps (2010), and Brennan et al. (1993) found no significant differences between lottery and control conditions. Mortagy et al. (1985) famously characterized their lottery condition as "a useless raffle," finding identical response rates (68.1\%) in both groups.

Most relevant to the present study, Denton et al. (1991) found lottery incentives ineffective among college alumni. This population shares key characteristics with Stanford alumni: high education, relative affluence, and an existing relationship with the survey sponsor. Zhang et al. (2017) provides a potential explanation, finding larger lottery effects among lower-income respondents. For affluent populations, the expected value of a modest lottery may fall below the threshold required to motivate behavior change.

Several studies have found that small prepaid cash incentives significantly outperform much larger lottery prizes. Larson and Chow (2003) found that a \$1 prepaid incentive outperformed a \$600 lottery, and Lorenzi et al. (1988) found lottery incentives performed significantly worse than \$1 prepaid cash among business executives. These findings suggest that the certainty of a small reward may be more motivating than the possibility of a large one, particularly for populations that do not need the money.

Recent meta-analytic evidence confirms the overall pattern. Göritz (2006) found an odds ratio of 1.19 for lottery incentives in web surveys, and Abdelazeem et al. (2023) reported a risk ratio of 1.12 for lotteries compared to 1.25 for cash. Lotteries work on average, but with smaller and more variable effects than guaranteed payments.

\subsection{Prize Structure and Framing Effects}

When researchers have a fixed incentive budget, they face a choice: offer many chances to win smaller prizes, or fewer chances to win larger prizes. The expected value can be held constant, but the psychological appeal may differ. Prospect theory suggests that people overweight small probabilities, which could favor either framing depending on whether respondents focus on the number of chances or the prize magnitude.

The evidence on prize structure is mixed. Deutskens et al. (2004) found that lotteries with smaller prizes and higher probability of winning were most effective. This finding aligns with the prediction that multiple chances feel psychologically better than a single long-shot opportunity.

Other studies find no significant differences across lottery amounts. Leung et al. (2002) and Halpern et al. (2011) found that varying the prize magnitude had no effect on response rates. Some evidence suggests larger prizes are more effective (Hubbard and Little 1988 found \$200 helped but \$50 to \$150 did not; Gajic et al. 2012 found high-value lotteries effective but low-value lotteries not), but this pattern is inconsistent.

Lottery type and timing also matter. Prepaid scratch tickets appear to outperform promised lottery entries (Olsen et al. 2012; Guo et al. 2016), possibly because their immediacy and tangibility increase perceived value. Tuten et al. (2004) found that immediate notification of lottery results outperformed delayed notification. Pre-incentive lotteries (unconditional entry) may be more effective than post-incentive lotteries (conditional on completion) (Sanchez-Fernandez et al. 2010).

Interactions with survey characteristics add further complexity. Spry et al. (1989) found lottery incentives helped only when combined with a short questionnaire, suggesting that incentive effectiveness depends on respondent burden. Cobanoglu and Cobanoglu (2003) found that lottery combined with other incentive types was effective, while lottery alone was not.

\subsection{The Understudied Question: Sample Composition}

Most studies of lottery incentives focus exclusively on response rates. A higher response rate is treated as a success. This framing ignores a critical question: do incentives change who responds?
If incentives motivate participation uniformly across the population, higher response rates translate directly into better data. If incentives disproportionately attract certain subgroups, they may increase response rates while simultaneously reducing representativeness. A survey with a higher response rate could have more nonresponse bias than one with a lower response rate, depending on which additional respondents the incentive recruits.

Several studies provide indirect evidence of differential effects. Laguilles et al. (2011) observed that males showed larger response increases for technology prizes while females responded more to dining-related prizes. Zhang et al. (2017) found larger lottery effects among low-income staff compared to higher-income staff. These patterns suggest incentives do not operate uniformly across demographic groups.

A few studies have examined composition directly. Robertson et al. (2005) found that a lottery incentive improved physician response rates and examined representativeness of responders versus nonresponders. Petrolia and Bhattacharjee (2009) found that incentives altered the education composition of their sample. Goyder et al. (1994) found that prepaid incentives yielded a more representative sample overall.

These studies raise the possibility that incentives could improve representativeness by drawing in groups that would otherwise be underrepresented. Lower-income respondents, for whom the expected value of a lottery represents a larger relative payoff, might be particularly motivated by financial incentives. If low-income individuals are underrepresented at baseline, incentives could reduce nonresponse bias.

The alternative possibility has received less empirical attention. Incentives might increase response most among groups already inclined to participate. In institutional populations like alumni, donors and active members may feel the reciprocity norm particularly strongly: they have given to the institution, and now the institution offers to give something back. If incentives disproportionately motivate the already-engaged, they would widen existing gaps in response propensity, increasing nonresponse bias despite increasing the overall response rate.

% ============================================
% METHODS
% ============================================

\section{Methods}

\subsection{Sample and Assignment}

The population was defined as people who received degrees from Stanford University between 1955 and 2007, who live in the continental USA or Canada, and for whom Stanford had an e-mail address. Stanford employees and their spouses, current and former University Trustees, current and former members of the Stanford Alumni Association Board of Directors, and those who had requested to receive no e-mail were excluded from the study.

Alumni were assigned to experimental conditions using systematic sampling. After removing those who had requested to receive no surveys (N = 2,499, 2,513, and 2,510 from the three groups), invitations were distributed sequentially in a 1-2-3-1-2-3 pattern, yielding three groups of alumni. Each individual was e-mailed a short invitation with a link to participate in the survey. The invitation was the same for all---except for the subject headers and the passages describing the incentive offered for their participation.

\begin{itemize}
    \item Sample 1 was offered no incentive to participate in the survey.
    \item Sample 2 was informed that five randomly selected respondents would each win a \$100 Visa gift card.
    \item Sample 3 was informed that one randomly selected respondent would win a \$500 Visa gift card.
\end{itemize}

The invitations were e-mailed on June 24, 2008, and three reminders were e-mailed to non-respondents at one-week intervals before the survey was taken off the web on July 20, 2008. A few invitations bounced back as undeliverable (N = 22, 29, and 24, respectively), yielding final participating samples of 2,430, 2,449, and 2,454. A balance check confirmed that the three experimental groups did not differ significantly on any measured demographic characteristic (all p > .16), validating the effectiveness of the systematic assignment procedure.

\subsection{Measures}

Names and demographic data came from Stanford's alumni database. The following characteristics served as covariates in the response rate models, capturing the dimensions along which response propensity varied in this population. Reference categories are noted in parentheses.

\textit{Sex.} Male (reference) or female.

\textit{Donor status.} Whether the alumnus had donated to Stanford since January 1, 2007 (reference: non-donor).

\textit{SAA membership.} Whether the alumnus was a member of the Stanford Alumni Association at the time of the survey (reference: non-member).

\textit{Degree level.} Alumni whose only degree from Stanford is a graduate degree (``graduate-only,'' reference) versus those who hold an undergraduate degree from Stanford (``undergraduate''). Alumni who have both undergraduate and graduate degrees from Stanford are classified as undergraduate alumni.

\textit{Graduation year cohort.} Graduation years grouped into five cohorts: 1955--1969 (reference), 1970--1979, 1980--1989, 1990--1999, and 2000--2007. Year of undergraduate graduation is recorded for alumni who have both undergraduate and graduate degrees.

\textit{Geographic region.} San Francisco Bay Area (reference), other active alumni regions, or non-active regions.

The response rate analyses (Analyses 1--3) draw on a dataset covering all 7,333 sample members, which includes treatment assignment, response status, and the demographic characteristics listed above. Age and race/ethnicity were not available in this file. Among the 2,342 respondents, a separate line-item file records responses to each survey item along with self-reported age and race/ethnicity. Age (D1) was recorded as a continuous measure in years (range: 21--88). Race/ethnicity (D2) was recorded using seven categories: White/Caucasian (reference), Asian/Pacific Islander/Indian, Latino/Hispanic, African-American/Black, Native American/American Indian, Mixed ethnicity/multiracial/other, and Prefer not to answer. These two variables are used as additional controls in the data quality analyses (Analysis 4) but are not available for the full sample.

\textit{Survey instrument.} The survey contained two batteries of attitudinal items. The Q1 battery comprised 26 items asking alumni to rate how much each activity made them feel connected to Stanford on a six-point scale (1 = a great deal, 2 = a lot, 3 = a moderate amount, 4 = only a little, 5 = not at all, 6 = not applicable). The Q2 battery comprised 16 items asking alumni to rate the extent to which each alumni service met their needs or interests (1 = fully, 2 = mostly, 3 = partly, 4 = only a little, 5 = not at all, 6 = not applicable). For response-quality analyses, ``not applicable'' responses (6) were treated as missing, as they represent inapplicability rather than a position on the attitude continuum.

\subsection{Analysis Plan}

We conducted four sets of analyses. The estimating equations and full statistical output for each are presented in the Analysis section; interpretive findings follow in the Results section.

\textit{Analysis 1: Overall incentive effect.} We compared response rates across the three experimental conditions using chi-square tests and one-sided tests (Fisher's exact and one-sided z-test) given our directional hypothesis. We then estimated bivariate and multivariate OLS regression models predicting survey response (0/1) as a function of treatment assignment and demographic covariates. We conducted joint Wald tests of the null hypothesis that $\beta_1 = 0$. The bivariate model takes the form:

\begin{equation}
\text{Responded}_i = \beta_0 + \beta_1 \cdot \text{AnyIncentive}_i + \epsilon_i
\end{equation}

The multivariate model adds the full covariate vector $\mathbf{X}_i$ described in Section 3.2:

\begin{equation}
\text{Responded}_i = \beta_0 + \beta_1 \cdot \text{AnyIncentive}_i + \mathbf{X}_i'\boldsymbol{\gamma} + \epsilon_i
\end{equation}

\textit{Analysis 2: Incentive framing.} To test whether lottery structure matters when expected value is held constant, we restricted the sample to the two incentive conditions (excluding the no-incentive control) and coded the treatment variable as 1 = one \$500 prize, 0 = five \$100 prizes (reference). We first compared raw response rates via chi-square test:

\begin{equation}
\text{Responded}_i = \beta_0 + \beta_1 \cdot \text{One500}_i + \epsilon_i
\end{equation}

We then estimated multivariate models adding $\mathbf{X}_i$, and tested for heterogeneous framing effects by adding treatment-by-covariate interaction terms with joint Wald tests:

\begin{equation}
\text{Responded}_i = \beta_0 + \beta_1 \cdot \text{One500}_i + \mathbf{X}_i'\boldsymbol{\gamma} + \epsilon_i
\end{equation}

\textit{Analysis 3: Subgroup effects.} We examined differential incentive effects across demographic subgroups through two complementary approaches: two-proportion z-tests comparing no-incentive and pooled-incentive response rates within each subgroup, and interaction models using the full sample with treatment-by-demographic terms and joint Wald tests.

\textit{Analysis 4: Data quality.} Among respondents (N = 2,342: 781 no-incentive, 1,561 incentive), we assessed response quality using five indicators: (1) item nonresponse---percentage of questions left unanswered per respondent, compared via t-test and multivariate OLS controlling for age and race/ethnicity; (2) straight-lining---identical responses across 10 or more Q1 items or 8 or more Q2 items, compared via chi-square test; (3) within-person response variance---standard deviation across answered items per battery, compared via t-test; (4) Mulligan non-differentiation scores---a pairwise-difference measure of response uniformity robust to scale endpoints \citep{anand2013}, estimated via t-test and multivariate OLS controlling for age and race/ethnicity; and (5) response distributions---chi-square tests on all 42 survey items comparing response distributions between conditions, with Bonferroni-Holm correction for multiple comparisons.


% ============================================
% ANALYSIS
% ============================================

\section{Analysis}

\subsection{Baseline Response Rates}

A total of 2,342 out of 7,333 alumni participated in the survey, yielding an overall response rate of 31.9\%. Response rates varied substantially across demographic and affiliation characteristics. Alumni with stronger institutional ties showed markedly higher participation: donors responded at 40.6\% compared to 26.5\% among non-donors ($\chi^2$ = 157.86, p < .0001); SAA members at 36.4\% versus 26.3\% among non-members ($\chi^2$ = 84.59, p < .0001); and undergraduate alumni at 35.0\% compared to 27.7\% among graduate-only alumni ($\chi^2$ = 43.79, p < .0001). Women responded at significantly higher rates (35.9\%) than men (29.5\%) ($\chi^2$ = 32.00, p < .0001). Among undergraduate alumni, response rates varied by graduation year cohort, with those who graduated in the 1970s showing the highest participation (38.3\%) and those in the 1980s and 1990s showing lower rates (32.6\% and 32.3\%, respectively) ($\chi^2$ = 11.38, p < .01).

\begin{figure}[H]
    \caption{Response rates by donor status, SAA membership, and degree level. Error bars represent 95\% confidence intervals.}
    \label{fig:response_affiliation}
    \centering
    \includegraphics[width=0.85\textwidth]{paper-figures/affiliation-response-rates.png}
\end{figure}

\begin{figure}[H]
    \caption{Response rates by sex and graduation year cohort (undergraduate alumni only). Error bars represent 95\% confidence intervals.}
    \label{fig:response_demographics}
    \centering
    \includegraphics[width=0.85\textwidth]{paper-figures/demographics-combined.png}
\end{figure}

These patterns are consistent across all three treatment conditions. As shown in Table~\ref{tab:subgroup_treatment_rr}, the demographic composition of respondents was similar regardless of incentive assignment, reinforcing that the baseline differences in response propensity reflect stable characteristics of the population rather than artifacts of the experimental conditions.

\begin{table}[H]
    \caption{Response rates by demographic subgroup and treatment condition.}
    \label{tab:subgroup_treatment_rr}
    \centering\small
    \input{paper-figures/subgroup-treatment-rr.tex}
\end{table}

\subsection{Analysis 1: Overall Incentive Effect on Response Rates}

Alumni offered any lottery incentive responded at 32.47\%, compared to 30.62\% in the no-incentive condition---a raw difference of 1.85 percentage points. A one-sided z-test yielded p = .055, a one-tailed Fisher's exact test p = .057, and a two-sided chi-square test p = .109.

In bivariate OLS models, the estimated incentive effect was 1.85 percentage points ($\beta$ = 0.019, SE = 0.011, p = .107). Adding demographic controls increased the point estimate slightly to 2.01 percentage points while reducing the p-value ($\beta$ = 0.020, SE = 0.011, p = .076). This pattern---where controls reduce the p-value rather than increase it---indicates that the demographic covariates are negatively correlated with treatment assignment, increasing statistical precision upon their inclusion.

\begin{table}[H]
    \caption{Bivariate OLS regression: incentive effect on survey response (0/1).}
    \label{tab:bivariate_results}
    \centering\small
    \begin{minipage}{0.5\textwidth}
        \input{paper-figures/ols-coef-bivariate.tex}
    \end{minipage}
\end{table}

The strongest predictors in the multivariate model were institutional affiliation measures. Donors were 12.0 percentage points more likely to respond than non-donors (p < .001); SAA members were 6.7 points more likely than non-members (p < .001); women were 5.8 points more likely than men (p < .001); and undergraduate alumni were 2.5 points more likely than graduate-only alumni (p = .033). Among graduation year cohorts, the 1990s cohort showed 4.9 points lower response than the 1955--1969 reference cohort (p = .004). All institutional affiliation effects substantially exceeded the incentive effect in magnitude.

\begin{table}[H]
    \caption{Multivariate OLS regression predicting survey response (0/1).}
    \label{tab:multivariate_results}
    \centering\small
    \input{paper-figures/ols-coef-multivariate.tex}
\end{table}

Joint Wald tests failed to reject the null hypothesis that both incentive coefficients equal zero (F = 1.60, p = .203). The 91 additional responses generated across the two incentive conditions came at a total prize cost of \$1,000, a marginal cost of \$10.99 per additional respondent.

\begin{figure}[H]
    \caption{Response rates by condition with 95\% confidence intervals. The pooled incentive rate combines both lottery framings.}
    \label{fig:response_by_condition}
    \centering
    \includegraphics[width=0.85\textwidth]{paper-figures/treatment-response-rates.png}
\end{figure}

\subsection{Analysis 2: Effect of Lottery Structure on Response Rates}

The two incentive conditions---five chances to win \$100 versus one chance to win \$500---had identical expected values (\$2.00 per respondent). Empirically, the two framings produced virtually identical response rates: 32.58\% for five \$100 prizes versus 32.36\% for one \$500 prize, a difference of 0.22 percentage points ($\chi^2$ = 0.03, p = .864).

Multivariate OLS models comparing the two incentive conditions directly confirmed this null result (F = 0.04, p = .836), with point estimates differing by less than 0.3 percentage points from the bivariate specification.

\begin{table}[H]
    \caption{Regression models comparing five \$100 prizes to one \$500 prize (no-incentive control excluded).}
    \label{tab:framing_comparison}
    \centering\small
    \input{paper-figures/ols-coef-framing.tex}
\end{table}

Interaction models testing heterogeneous framing effects across demographic subgroups found no individual interaction to be statistically significant, and the joint Wald test was far from significant (F = 1.03, p = .414).

\begin{figure}[H]
    \caption{Response rates by incentive framing with 95\% confidence intervals. The two structures produced identical participation despite different probability-prize combinations.}
    \label{fig:framing_rrs}
    \centering
    \includegraphics[width=0.85\textwidth]{paper-figures/framing-comparison.png}
\end{figure}

\subsection{Analysis 3: Heterogeneous Treatment Effects Across Demographic Subgroups}

Two-proportion z-tests comparing no-incentive and pooled-incentive response rates within each subgroup revealed a consistent pattern: incentive effects were largest among groups with the highest baseline response rates. Women showed a 3.4 percentage point increase ($\chi^2$ = 3.15, p = .076), while men showed only 0.9 points ($\chi^2$ = 0.38, p = .539). Donors increased by 3.4 points ($\chi^2$ = 2.96, p = .085), while non-donors increased by 1.0 point ($\chi^2$ = 0.46, p = .498).

\begin{table}[H]
    \caption{Two-proportion z-tests comparing response rates between no-incentive and pooled-incentive conditions within demographic subgroups.}
    \label{tab:subgroup_effects}
    \centering\small
    \begin{minipage}{0.45\textwidth}
        \input{paper-figures/subgroup-ztests.tex}
    \end{minipage}
\end{table}

Formal interaction models including treatment-by-demographic terms found no individual interaction to be statistically significant. The largest point estimates---for donors ($\hat{\gamma}$ = 0.031, p = .208) and women ($\hat{\gamma}$ = 0.027, p = .267)---were directionally consistent with the z-test patterns but well below conventional thresholds. Joint Wald tests failed to reject homogeneous treatment effects (F = 0.67, p = .755).

\subsection{Analysis 4: Incentive Effects on Response Quality}

\subsubsection{Item Nonresponse}

Respondents in the incentive condition left significantly fewer questions unanswered (1.99\%) compared to the no-incentive condition (2.81\%), a difference of 0.82 percentage points (t = 2.089, p = .037). This represents a 29\% reduction in item nonresponse. In multivariate OLS controlling for age and race/ethnicity, the incentive coefficient was $\beta$ = $-$0.823, SE = 0.393, p = .037 (N = 2,342; $R^2$ = 0.014). The effect persisted after accounting for demographic composition.

Five items showed particularly high nonresponse across both conditions: Q2O (8.07\%), Q2P (5.47\%), Q2N (5.34\%), Q2L (4.91\%), and Q2G (4.82\%)---all from the Q2 battery, all including a ``Not Applicable'' response option.

\subsubsection{Satisficing Behavior}

Straight-lining was virtually nonexistent in both conditions. For the Q1 battery, only 3 of 2,342 respondents (0.13\%) gave identical responses to 10 or more items, with identical rates across conditions ($\chi^2$ = 0.00, p = 1.000). The Q2 battery showed slightly higher but still minimal rates: 1.47\% in no-incentive versus 0.63\% in incentive conditions ($\chi^2$ = 2.54, p = .111).

Within-person response variance showed no significant differences between conditions. Mean standard deviations were nearly identical for Q1 (1.51 vs.\ 1.49, t = $-$0.88, p = .380) and Q2 (1.67 vs.\ 1.64, t = $-$1.25, p = .211). Mulligan non-differentiation scores---a pairwise-difference measure more robust to scale endpoints than simple variance---corroborated this result: mean ND scores were virtually identical between conditions for Q1, Q2, and all 42 items combined, with no bivariate or multivariate test approaching significance.

\subsubsection{Response Distributions}

Chi-square tests on all 42 survey items with Bonferroni-Holm correction found no item with a statistically significant distributional difference between conditions after correction. The two items with the smallest unadjusted p-values---Q1R and Q1A, both p = .004---had adjusted p-values of .146.

\begin{table}[H]
    \caption{Chi-square tests for response distribution differences between conditions. Top 10 items by unadjusted $p$-value shown. None remained significant after Bonferroni-Holm correction.}
    \label{tab:response_distributions}
    \centering\small
    \begin{minipage}{0.65\textwidth}
        \input{paper-figures/response-distribution-tests.tex}
    \end{minipage}
\end{table}



% ============================================
% RESULTS
% ============================================

\section{Results}

\subsection{Alumni with Stronger Institutional Ties Are Substantially More Likely to Respond}

A total of 2,342 out of 7,333 alumni participated in the survey, yielding an overall response rate of 31.9\%.\footnote{People who had asked to receive no surveys were removed from the three random samples (N = 2,499, 2,513, and 2,510), leaving samples of 2,452, 2,478, and 2,478. Invitations were e-mailed to these people, and a few bounced back as undeliverable (N = 22, 29, and 24, respectively), yielding final participating samples of 2,430, 2,449, and 2,454.}

Before examining the experimental effects, it is important to establish who responds to this survey without any incentive---because this baseline distribution shapes what any incentive can and cannot accomplish. Response rates varied substantially by institutional connection: donors responded at 40.6\% versus 26.5\% among non-donors; SAA members at 36.4\% versus 26.3\% among non-members; undergraduate alumni at 35.0\% versus 27.7\% among graduate-only alumni. Women responded at markedly higher rates (35.9\%) than men (29.5\%). All of these differences were highly significant (all p < .0001).

These differences were consistent across all measured characteristics: donors, SAA members, women, and undergraduate alumni all participated at substantially higher rates than their counterparts. This baseline heterogeneity is not merely descriptive context---it shapes what any incentive can and cannot accomplish.

\subsection{Lottery Incentives Produce a Small, Inconsistent Increase in Response Rates}

Lottery incentives increased overall response rates by approximately 2 percentage points (30.6\% to 32.5\%), but this effect was small and statistically marginal. The most permissive test---a one-sided z-test reflecting our directional hypothesis---yielded p = .055, just above the conventional threshold. Two-sided tests and joint Wald tests failed to reject the null of no effect.

The incentive effect's sensitivity to the inclusion of demographic controls is itself informative. In bivariate models the effect is p = .107; once demographic controls are added it moves to p = .076. Neither crosses the .05 threshold. Statistical power to detect a 2 percentage point effect is inherently limited in a study of this size, and we cannot rule out that a true but small effect exists. What we can say is that the evidence for a meaningful boost in participation is weak---and that the 91 additional respondents generated cost \$10.99 each in prize expenditure. Whether that cost is justified depends critically on which respondents the incentives recruited, a question we take up next.

\subsection{Lottery Structure Has No Effect: Framing Does Not Matter When Expected Value Is Equal}

The two lottery structures---five chances at \$100 versus one chance at \$500---produced virtually identical response rates (32.58\% and 32.36\%), a gap of 0.22 percentage points that was nowhere near statistical significance. This null result held across all demographic subgroups as well.

This finding has a clear practical implication for survey researchers working with constrained budgets: when designing a lottery incentive for a similar population, the structure of the prize is irrelevant. A single large prize and multiple smaller prizes of equal total value are interchangeable. Respondents in this population appear to evaluate lotteries on the basis of expected value rather than through the psychological biases---overweighting of probabilities, sensitivity to prize magnitude---that prospect theory might predict.

\subsection{Incentives Widen Existing Participation Gaps Rather Than Narrowing Them}

The most consequential finding of this study is not the overall response rate increase but \textit{who} generated it. Incentives increased participation most among the groups already most likely to respond: women (+3.4 pp) versus men (+0.9 pp), and donors (+3.4 pp) versus non-donors (+1.0 pp).

The practical implication is that incentives did not broaden the survey's reach---they deepened participation among groups with high baseline response propensity. The response rate difference between women and men widened from 4.7 to 7.3 percentage points, and between donors and non-donors from 12.5 to 14.9 percentage points. Though these differential effects did not achieve conventional statistical significance, the consistent pattern across subgroups suggests that incentives amplified rather than reduced existing differences in who was likely to respond.

\subsection{Data Quality Is Preserved: Incentives Improve Completion Without Biasing Responses}

A concern sometimes raised about financial incentives is that they attract inattentive respondents who rush through surveys and provide low-quality answers. The evidence here runs counter to that expectation. Incentivized respondents left \textit{fewer} questions unanswered (1.99\% vs.\ 2.81\%), a 29\% reduction in item nonresponse that persisted after controlling for demographics. The incentive was associated with more thorough completion of the survey, not less.

Beyond item nonresponse, every other quality indicator showed no significant difference by condition. Straight-lining was negligible and equal across groups. Within-person response variance was nearly identical. Response distributions on all 42 substantive items showed no significant differences after multiple-comparison correction. Incentivized and non-incentivized alumni evaluated their Stanford connection and satisfaction with alumni services in statistically indistinguishable ways.

This is an important reassurance: whatever compositional effects incentives introduced, they did not corrupt the substantive answers of those who responded. The concern is not that incentivized respondents gave worse data---it is that incentives recruited more of the same kinds of respondents who would have participated anyway.


% ============================================
% DISCUSSION
% ============================================

\section{Discussion}

\subsection{Summary of Main Findings}

Three principal findings emerged from this experiment. First, lottery incentives produced a small and inconsistent increase in response rates---approximately 2 percentage points---that was sensitive to the inclusion of demographic controls (p-values ranging from .055 to .107) and failed to achieve significance under joint Wald tests (p $\approx$ .20). The 91 additional responses generated came at a total prize cost of \$1,000, a marginal cost of \$10.99 per additional respondent. Second, incentive framing had no detectable effect when expected value was held constant: five chances at \$100 produced virtually identical response rates as one chance at \$500 (difference = 0.22 percentage points, p = .864). Third, incentives disproportionately increased response among subgroups already over-represented at baseline---particularly women and donors---amplifying rather than reducing existing differences in who was likely to respond.

\subsection{Theoretical Implications}

\subsubsection{Population Characteristics and Lottery Effectiveness}

The weak incentive effect is consistent with prior evidence that lotteries perform poorly in affluent, highly educated, and institutionally connected populations (Cottrell et al.\ 2015; Denton et al.\ 1991). For Stanford alumni, the \$2 expected value likely falls well below any motivational threshold. The pre-existing relationship between alumni and the sponsor may further reduce the marginal impact of a financial offer, as institutional affiliation already provides substantial leverage for many potential respondents.

\subsubsection{Prize Structure and Framing Effects}

The null framing effect suggests that this population evaluated the two lotteries on the basis of expected value rather than through the heuristic processing that prospect theory might predict. The explicit disclosure of prize amounts and number of winners likely facilitated rational calculation of expected value, and a highly educated population may be particularly capable of doing so. This contrasts with Deutskens et al.\ (2004), who found higher-probability, smaller-prize lotteries more effective---a discrepancy that may reflect differences in population affluence or the absolute magnitude of the prizes offered. Practically, researchers designing lottery incentives for similar populations need not devote resources to optimizing the probability-prize combination.

\subsubsection{Respondent Composition and Nonresponse Bias}

The differential effects across subgroups speak directly to the understudied question of whether lottery incentives change \textit{who} responds, not just \textit{how many} respond. Rather than broadening participation among underrepresented alumni, incentives concentrated additional responses among those already most likely to participate. Two theoretical accounts are consistent with this pattern. Leverage-saliency theory (Groves, Singer, \& Corning 2000) suggests that the monetary offer was simply more salient to respondents already engaged with the institution. The norm of reciprocity offers a complementary explanation: donors, having already given to Stanford, may feel the incentive activates an obligation of reciprocal exchange that non-donors do not share. Either way, the consistent directional pattern suggests that incentives operated through mechanisms tied to existing institutional engagement rather than compensating for its absence.

\subsection{Methodological Implications}

These findings suggest that lottery incentives warrant careful consideration in institutional survey contexts. Cost per additional response is an incomplete metric when incentives disproportionately attract respondents who are already over-represented: the cost per additional representative response may substantially exceed \$10.99. Researchers should assess baseline heterogeneity in response propensity before deploying incentives, and consider whether alternative strategies---targeted outreach to underrepresented groups, prepaid incentives (Larson \& Chow 2003), or additional contact attempts---would better serve representativeness goals without differential effects. Because all conditions in this study received identical follow-up protocols, we cannot assess whether incentives reduce the optimal number of reminders needed---a potential cost offset that warrants investigation in future designs that manipulate both incentives and follow-up intensity (Olsen et al.\ 2012). Evaluations of incentive effectiveness should routinely examine the demographic breakdown of respondents across conditions, not merely overall response rates.

\subsection{Limitations and Future Directions}


\subsection{Conclusion}

Lottery incentives produced a small and inconsistent increase in response rates while disproportionately motivating subgroups already over-represented among respondents, widening existing gaps in participation. Data quality was preserved---incentivized respondents completed the survey more thoroughly and answered substantive questions in statistically indistinguishable ways---but this reassurance does not address the compositional concern. When incentives differentially attract high-propensity respondents, they can increase nonresponse bias even as they increase total response.

For researchers designing surveys within institutional populations---alumni, employees, members---the practical recommendation is clear: before deploying a lottery incentive, assess who is already likely to respond and whether the incentive will reach those who are not. If baseline response propensity is heterogeneous and strongly correlated with the survey's substantive variables, a lottery incentive may deliver a marginal increase in total response while widening existing differences in who participates. The appropriate benchmark for evaluating incentive strategies is not response rate alone, but whether the additional responses improve the representativeness and analytical value of the data.

\end{document}